% Topic 1.1.11 — Constrained Extrema. Lagrange Multipliers.

\section{Условный экстремум. Метод Лагранжа}
\label{sec:1.1.11-conditional-extrema}

\begin{ticket}
Поиск экстремумов функций от многих переменных. Условный экстремум функции нескольких переменных. Метод множителей Лагранжа, необходимые и достаточные условия условного экстремума.
\end{ticket}

\subsection*{Теория}

Рассматриваем вещественнозначные функции $f \colon U \to \R$ на
открытом множестве $U \subseteq \R^n$, обладающие достаточной
гладкостью для всех используемых ниже частных производных первого и
второго порядков.

\subsubsection*{Безусловный экстремум}

\begin{definition}
\label{def:hessian}
Если~$f$ имеет в точке $\vect{x} \in U$ все частные производные
второго порядка, её \emph{гессианом} называется симметричная матрица
\[
  H_f(\vect{x}) \;=\; \Bigl(\tfrac{\partial^2 f}{\partial x_i \partial x_j}(\vect{x})\Bigr)_{i,j = 1}^{n}.
\]
(Симметричность опирается на теорему Шварца о равенстве смешанных
производных при их непрерывности.)
\end{definition}

\begin{theorem}[Необходимое условие внутреннего экстремума]
\label{thm:necessary-uncon}
Если~$f$ дифференцируема во внутренней точке~$\vect{x}_*$
множества~$U$ и имеет в ней локальный экстремум, то
$\nabla f(\vect{x}_*) = \vect{0}$. Точка с
$\nabla f(\vect{x}_*) = \vect{0}$ называется \emph{критической точкой}
функции~$f$.
\end{theorem}
\proofof{necessary-uncon}

\begin{theorem}[Достаточное условие через гессиан]
\label{thm:sufficient-uncon}
Пусть~$f$ дважды непрерывно дифференцируема в окрестности
точки~$\vect{x}_*$ и $\nabla f(\vect{x}_*) = \vect{0}$.
\begin{enumerate}[label=(\alph*)]
  \item Если $H_f(\vect{x}_*)$ положительно определена, то $\vect{x}_*$~---
        строгий локальный минимум.
  \item Если $H_f(\vect{x}_*)$ отрицательно определена, то $\vect{x}_*$~---
        строгий локальный максимум.
  \item Если $H_f(\vect{x}_*)$ неопределена (имеет собственные
        значения обоих знаков), то $\vect{x}_*$~--- седловая точка,
        не являющаяся локальным экстремумом.
\end{enumerate}
Граничные случаи (полуопределённый, но не определённый гессиан)
данным признаком не разрешаются.
\end{theorem}
\proofof{sufficient-uncon}

Определённость $H_f(\vect{x}_*)$ удобно проверять по критерию
Сильвестра (\cref{thm:sylvester}).

\subsubsection*{Условный экстремум}

Перейдём к задаче нахождения экстремума~$f$ при ограничениях
\[
  g_1(\vect{x}) = 0, \;\dots, \; g_m(\vect{x}) = 0 \qquad (m < n),
\]
где каждая $g_i \colon U \to \R$ непрерывно дифференцируема.

\begin{definition}
\label{def:constrained-extremum}
Точка $\vect{x}_* \in U$ называется \emph{условным локальным
минимумом} (\emph{максимумом}) функции~$f$ при заданных ограничениях,
если $g_i(\vect{x}_*) = 0$ для всех~$i$ и существует окрестность
$V \ni \vect{x}_*$, в которой $f(\vect{x}_*) \le f(\vect{x})$
(соответственно $\ge$) для всякой $\vect{x} \in V$, удовлетворяющей
$g_i(\vect{x}) = 0$ при всех~$i$.
\end{definition}

\begin{definition}
\label{def:regularity}
Точка~$\vect{x}_*$ на множестве ограничений (то есть с
$g_i(\vect{x}_*) = 0$ при всех~$i$) называется \emph{регулярной},
если градиенты $\nabla g_1(\vect{x}_*), \dots, \nabla g_m(\vect{x}_*)$
линейно независимы.
\end{definition}

\begin{definition}
\label{def:lagrangian}
\emph{Лагранжианом} задачи на условный экстремум называется
\[
  L(\vect{x}, \vect{\lambda}) \;=\; f(\vect{x}) - \sum_{i=1}^{m} \lambda_i\, g_i(\vect{x}),
  \qquad \vect{\lambda} = (\lambda_1, \dots, \lambda_m) \in \R^m;
\]
скаляры $\lambda_i$ называются \emph{множителями Лагранжа}.
\end{definition}

\begin{theorem}[Необходимое условие Лагранжа]
\label{thm:lagrange-necessary}
Пусть~$\vect{x}_*$~--- регулярная точка условного локального
экстремума функции~$f$. Тогда существуют единственные множители
$\lambda_1^*, \dots, \lambda_m^*$, для которых
\[
  \nabla f(\vect{x}_*) \;=\; \sum_{i=1}^{m} \lambda_i^*\, \nabla g_i(\vect{x}_*).
\]
Эквивалентно, $(\vect{x}_*, \vect{\lambda}^*)$~--- критическая точка
лагранжиана: $\nabla_{\vect{x}} L(\vect{x}_*, \vect{\lambda}^*) = \vect{0}$
вместе с равенствами $g_i(\vect{x}_*) = 0$.
\end{theorem}
\proofof{lagrange-necessary}

\begin{theorem}[Достаточное условие условного экстремума]
\label{thm:lagrange-sufficient}
Пусть $\vect{x}_*$~--- регулярная точка с множителями
$\vect{\lambda}^*$, удовлетворяющая
$\nabla_{\vect{x}} L(\vect{x}_*, \vect{\lambda}^*) = \vect{0}$ и
$g_i(\vect{x}_*) = 0$. Положим $T = \{\vect{h} \in \R^n : \langle \nabla g_i(\vect{x}_*), \vect{h} \rangle = 0
\text{ для всех } i\}$~--- касательное пространство многообразия
ограничений в точке~$\vect{x}_*$, и пусть
$H = H_{L(\cdot, \vect{\lambda}^*)}(\vect{x}_*)$~--- гессиан функции
$L(\,\cdot\,, \vect{\lambda}^*)$ в точке~$\vect{x}_*$.
\begin{enumerate}[label=(\alph*)]
  \item Если $\vect{h}\T H \vect{h} > 0$ для всякого ненулевого
        $\vect{h} \in T$, то $\vect{x}_*$~--- строгий условный
        локальный минимум.
  \item Если $\vect{h}\T H \vect{h} < 0$ для всякого ненулевого
        $\vect{h} \in T$, то $\vect{x}_*$~--- строгий условный
        локальный максимум.
  \item Если $H$, ограниченная на~$T$, неопределена, то $\vect{x}_*$
        не является условным экстремумом.
\end{enumerate}
\end{theorem}
\proofof{lagrange-sufficient}

\begin{remark}
На практике решают систему из $n + m$ уравнений
$\nabla_{\vect{x}} L = \vect{0}$, $\nabla_{\vect{\lambda}} L = \vect{0}$
(второе~--- это сами ограничения) относительно $n + m$ неизвестных
$(\vect{x}, \vect{\lambda})$, а затем классифицируют каждое решение,
сужая $H_L$ на касательное пространство~$T$. Взгляд, при котором
безусловный экстремум обобщается на условный заменой $f$ на $L$,~---
дверь в выпуклую оптимизацию (условия ККТ, двойственность); см.\
\cref{sec:2.4-convex-optimization}; необходимые предварительные
сведения из исчисления функций многих переменных~--- в
\cite{Zorich1981}.
\end{remark}

\subsection*{Доказательства}

\begin{proof}[\Cref{thm:necessary-uncon}]
\proofanchor{necessary-uncon}
Пусть~$\vect{x}_*$~--- локальный минимум (максимум зеркален). Для
каждого~$i$ функция одной переменной
$\varphi_i(t) = f(\vect{x}_* + t\vect{e}_i)$ дифференцируема в $t = 0$
и $\varphi_i'(0) = \partial f / \partial x_i (\vect{x}_*)$
(\cref{def:partial-derivative}); кроме того, она имеет локальный
минимум в $t = 0$. По теореме Ферма (\cref{thm:fermat})
$\varphi_i'(0) = 0$, то есть
$\partial f / \partial x_i (\vect{x}_*) = 0$ при всех~$i$~--- значит,
$\nabla f(\vect{x}_*) = \vect{0}$.
\end{proof}

\begin{proof}[\Cref{thm:sufficient-uncon}]
\proofanchor{sufficient-uncon}
Разложение Тейлора~$f$ в точке~$\vect{x}_*$ до второго порядка
(\cref{thm:taylor-peano}, обобщённое на несколько переменных,
применяемое вдоль прямой $t \mapsto \vect{x}_* + t \vect{h}$) с
учётом $\nabla f(\vect{x}_*) = \vect{0}$ даёт
\[
  f(\vect{x}_* + \vect{h}) - f(\vect{x}_*) \;=\; \tfrac{1}{2}\, \vect{h}\T H_f(\vect{x}_*)\, \vect{h} + o(\|\vect{h}\|^2).
\]
\emph{(а) Положительно определённый гессиан.} Пусть $\mu > 0$~---
наименьшее собственное значение $H_f(\vect{x}_*)$ (вещественное и
положительное по спектральной теореме, \cref{thm:spectral}). Тогда
$\vect{h}\T H_f(\vect{x}_*) \vect{h} \ge \mu \|\vect{h}\|^2$ при
всех~$\vect{h}$, а член $o(\|\vect{h}\|^2)$ при достаточно малых
$\|\vect{h}\|$ не превосходит $\varepsilon \|\vect{h}\|^2$. Выбирая
$\varepsilon < \mu/2$, получаем
$f(\vect{x}_* + \vect{h}) - f(\vect{x}_*) \ge (\mu/2 - \varepsilon)\|\vect{h}\|^2 > 0$
для всякого ненулевого малого~$\vect{h}$. Значит, $\vect{x}_*$~---
строгий локальный минимум.

\emph{(б) Отрицательно определённый.} Применим (а) к~$-f$.

\emph{(в) Неопределённый.} Возьмём собственные векторы
$\vect{u}, \vect{v}$ матрицы $H_f(\vect{x}_*)$ с собственными
значениями разных знаков. Из приведённого разложения видно, что вдоль
прямой через~$\vect{x}_*$ в направлении~$\vect{u}$ функция~$f$
локально строго возрастает, а в направлении~$\vect{v}$~--- строго
убывает. Оба поведения исключают наличие локального экстремума.
\end{proof}

\begin{proof}[\Cref{thm:lagrange-necessary}]
\proofanchor{lagrange-necessary}
Пусть~$\vect{x}_*$~--- регулярная точка условного экстремума и
$M = \{\vect{x} \in U : g_1(\vect{x}) = \dots = g_m(\vect{x}) = 0\}$~---
многообразие ограничений. Регулярность (линейная независимость
$\nabla g_1(\vect{x}_*), \dots, \nabla g_m(\vect{x}_*)$) по теореме о
неявной функции (стандартный факт из~\cite{Zorich1981}) влечёт, что
вблизи~$\vect{x}_*$ множество~$M$~--- гладкое многообразие
размерности $n - m$, причём его касательное пространство в~$\vect{x}_*$
есть
\[
  T_{\vect{x}_*} M \;=\; \{\vect{v} \in \R^n : \langle \nabla g_i(\vect{x}_*), \vect{v} \rangle = 0 \text{ при всех } i\}.
\]

Для произвольного касательного вектора $\vect{v} \in T_{\vect{x}_*} M$
возьмём гладкую кривую
$\gamma \colon (-\varepsilon, \varepsilon) \to M$, для которой
$\gamma(0) = \vect{x}_*$ и $\gamma'(0) = \vect{v}$. Композиция
$f \circ \gamma$ имеет локальный экстремум в $t = 0$, поэтому по
\cref{thm:fermat}
$(f \circ \gamma)'(0) = \langle \nabla f(\vect{x}_*), \vect{v} \rangle = 0$.
Это верно для каждого $\vect{v} \in T_{\vect{x}_*} M$, поэтому
\[
  \nabla f(\vect{x}_*) \;\in\; (T_{\vect{x}_*} M)^{\perp}
  \;=\; \spn(\nabla g_1(\vect{x}_*), \dots, \nabla g_m(\vect{x}_*));
\]
второе равенство имеет место, поскольку ортогональным дополнением
подпространства, заданного $m$ независимыми линейными функционалами,
служит линейная оболочка этих функционалов. Значит,
$\nabla f(\vect{x}_*) = \sum \lambda_i^* \nabla g_i(\vect{x}_*)$ при
единственных скалярах~$\lambda_i^*$ (единственность~--- из линейной
независимости градиентов $\nabla g_i$).
\end{proof}

\begin{proof}[\Cref{thm:lagrange-sufficient}]
\proofanchor{lagrange-sufficient}
Доказательство сводит достаточный признак безусловного экстремума
(\cref{thm:sufficient-uncon}) к ограничению на многообразие
ограничений~$M$, которое вблизи~$\vect{x}_*$ есть гладкое
подмногообразие размерности $n - m$. Введём локальные координаты
$\vect{u} \in \R^{n - m}$ на~$M$ с центром в~$\vect{x}_*$ через
параметризацию $\vect{x} = \Phi(\vect{u})$, доставляемую теоремой о
неявной функции. Композиция
$\widetilde{f}(\vect{u}) = f(\Phi(\vect{u}))$ имеет локальный
экстремум в $\vect{u} = \vect{0}$, и стандартное вычисление даёт
\[
  \widetilde{f}(\vect{u}) - \widetilde{f}(\vect{0}) \;=\; \tfrac{1}{2}\, \vect{h}(\vect{u})\T H \vect{h}(\vect{u}) + o(\|\vect{u}\|^2),
\]
где $\vect{h}(\vect{u}) = \Phi'(\vect{0})\,\vect{u} + o(\|\vect{u}\|)$
параметризует касательное пространство~$T = T_{\vect{x}_*} M$ вблизи
нуля. Определённость~$H$ на~$T$ (пункты (а), (б), (в)) задаёт знак
$\widetilde{f}(\vect{u}) - \widetilde{f}(\vect{0})$ вблизи
$\vect{u} = \vect{0}$~--- в точности так же, как в
\cref{thm:sufficient-uncon}. Подробное вычисление см.\ в
\cite{Zorich1981}.
\end{proof}

\subsection*{Примеры}

\begin{example}[Безусловный: положительно определённая квадратичная форма]
\label{ex:uncon-quadratic}
Для $f(x, y) = x^2 + y^2 - 2x + 4y$ градиент равен
$\nabla f(x, y) = (2x - 2, 2y + 4)$. Приравнивая к нулю, получаем
единственную критическую точку $\vect{x}_* = (1, -2)$. Гессиан~---
постоянная матрица $H_f = \begin{pmatrix} 2 & 0 \\ 0 & 2 \end{pmatrix} = 2 \mat{I}$,
положительно определённая (по Сильвестру: $M_1 = 2 > 0$, $M_2 = 4 > 0$).
По \cref{thm:sufficient-uncon}\,(а) точка $\vect{x}_*$~--- строгий
локальный минимум, причём $f(\vect{x}_*) = 1 + 4 - 2 - 8 = -5$.
Поскольку~$f$~--- положительно определённая квадратичная форма, этот
локальный минимум одновременно является единственным глобальным.
\end{example}

\begin{example}[Метод Лагранжа на окружности]
\label{ex:lagrange-circle}
Найдём экстремумы $f(x, y) = x + y$ при условии
$g(x, y) = x^2 + y^2 - 1 = 0$. Лагранжиан:
$L = (x + y) - \lambda(x^2 + y^2 - 1)$, и
\[
  \nabla_{\vect{x}} L = (1 - 2\lambda x,\; 1 - 2\lambda y),
\]
откуда $1 = 2\lambda x = 2 \lambda y$, то есть $x = y$. Совмещая с
$x^2 + y^2 = 1$, получаем двух кандидатов:
$(1/\sqrt{2}, 1/\sqrt{2})$ при $f = \sqrt{2}$ и
$(-1/\sqrt{2}, -1/\sqrt{2})$ при $f = -\sqrt{2}$. Обе точки регулярны
($\nabla g = (2x, 2y) \neq \vect{0}$ на единичной окружности).
Окружность компактна, а~$f$ непрерывна, поэтому максимум $\sqrt{2}$
достигается в $(1/\sqrt{2}, 1/\sqrt{2})$, а минимум $-\sqrt{2}$~---
в $(-1/\sqrt{2}, -1/\sqrt{2})$.
\end{example}

\begin{example}[Метод Лагранжа с линейным ограничением]
\label{ex:lagrange-linear}
Найдём максимум $f(x, y) = xy$ при условии $g(x, y) = x + y - 4 = 0$.
$L = xy - \lambda(x + y - 4)$,
$\nabla_{\vect{x}} L = (y - \lambda, x - \lambda)$, откуда
$x = y = \lambda$, и из ограничения $x = y = 2$. Максимальное
значение $f(2, 2) = 4$. (Геометрическая проверка: по неравенству о
среднем арифметическом и геометрическом
$xy \le ((x + y)/2)^2 = 4$ на множестве ограничения, причём равенство
достигается лишь при $x = y$.)
\end{example}

\begin{problem}
\label{prob:rectangle-ellipse}
Найдите прямоугольник наибольшей площади, вписанный в эллипс
$\frac{x^2}{a^2} + \frac{y^2}{b^2} = 1$ так, чтобы его стороны были
параллельны осям координат ($a, b > 0$).
\end{problem}

\begin{proof}[Решение]
По симметрии прямоугольник имеет вершины $(\pm x, \pm y)$ с $x, y > 0$
на эллипсе. Его площадь равна $A = 4xy$. Максимизируем
$f(x, y) = 4xy$ при условии
$g(x, y) = x^2/a^2 + y^2/b^2 - 1 = 0$.

Лагранжиан: $L = 4xy - \lambda(x^2/a^2 + y^2/b^2 - 1)$, и
\[
  \frac{\partial L}{\partial x} = 4y - \frac{2\lambda x}{a^2} = 0,
  \qquad
  \frac{\partial L}{\partial y} = 4x - \frac{2\lambda y}{b^2} = 0.
\]
Из первого уравнения $\lambda = 2 a^2 y / x$; подставляя во второе,
$4x = (2/b^2)(2 a^2 y / x) y$, то есть $x^2 b^2 = a^2 y^2$, и
$y = (b/a) x$. Подставляя в ограничение,
$x^2/a^2 + (b/a)^2 x^2 / b^2 = 2 x^2 / a^2 = 1$, откуда
$x = a/\sqrt{2}$, $y = b/\sqrt{2}$, а максимальная площадь равна
$A = 4 (a/\sqrt{2})(b/\sqrt{2}) = 2 a b$.

Точка регулярна (градиент ограничения
$\nabla g = (2x/a^2, 2y/b^2)$ ненулевой на эллипсе). Замкнутый эллипс
(пересечённый с первым квадрантом для нашей параметризации) компактен,
$f$ непрерывна, поэтому максимум достигается, и единственная найденная
стационарная точка его и реализует.
\end{proof}
